
%(BEGIN_QUESTION)
% Copyright 2006, Tony R. Kuphaldt, released under the Creative Commons Attribution License (v 1.0)
% This means you may do almost anything with this work of mine, so long as you give me proper credit

Mange væsker gjennomgår prosesser der molekylene splittes i positivt og negativt ladede ioner. Flytende {\it ioniske} forbindelser splittes helt eller nesten helt i ioner, mens bare en liten andel molekyler i flytende {\it kovalente} forbindelser splittes i ioner. Prosessen kalles {\it dissosiasjon} for ioniske forbindelser, og {\it ionisering} for kovalente forbindelser.

Smeltet salt (NaCl) er et eksempel på førstnevnte, mens rent vann er et eksempel på sistnevnte. Høy ionekonsentrasjon i smeltet salt forklarer god ledningsevne, mens lav ionekonsentrasjon i rent vann forklarer at det ofte regnes som isolator.

Selv om ioniseringen i vann er liten, er den ikke null. Undersøk molariteten (mol per liter, $M$) for både anioner og kationer i rent vann ved 25$^{o}$ C.

\vskip 10pt

Hint: gode kilder er lærestoff om {\it pH} og {\it pH-måling}.

\underbar{file i00613}
%(END_QUESTION)





%(BEGIN_ANSWER)

\noindent
I rent vann ved 25$^{o}$ C:

\vskip 10pt

[H$^{+}$] = 1.00 $\times$ $10^{-7}$ $M$

\vskip 10pt

[OH$^{-}$] = 1.00 $\times$ $10^{-7}$ $M$

\vskip 10pt

For rent vann og fortynnede vannløsninger ved 25$^{o}$ C er produktet av disse to molaritetene (nær) lik 1.00 $\times$ $10^{-14}$. Dette er {\it ioniseringskonstanten for vann}, $K_W$:

$$K_W = [\hbox{H}^{+}] [\hbox{OH}^{-}] = 1.00 \times 10^{-14}$$

%(END_ANSWER)





%(BEGIN_NOTES)

Elektrisk ledningsevne i en væske er i praksis en definisjonsnær test på om stoffet er ionisk eller kovalent (molekylært).

%INDEX% Kjemi, ion: dissosiasjon (definert)
%INDEX% Kjemi, ion: ionisering (definert)

%(END_NOTES)

